% Reader question: when may a sender continue relative to receiver readiness?
% Claim: buffered sends and synchronous rendezvous have different blocking rules.
% Evidence: two schematic schedules on a shared ordinal time axis.
% Distinguish: nonblocking send, receiver progress, joint event and optional timeout.
% Grammar: paired swimlanes with aligned time and a keyed semantics table.
% Acceptance: find the send that proceeds before its receiver is ready.
\begin{minipage}{11.1cm}
\centering\pdfigurehue{pdcobalt}%
\begin{tikzpicture}[pd figure,x=1cm,y=1cm]
  \def\xzero{2.05}\def\xe{9.85}\def\dy{1.2}\def\py{-4.2}
  \def\sendx{3.9}\def\recvx{7.8}
  \node[pd title,anchor=west] at (0,1.15) {Asynchronous actor send};
  \node[pd label,anchor=east,text width=1.75cm,align=right] at (1.75,0) {sender actor};
  \node[pd label,anchor=east,text width=1.75cm,align=right] at (1.75,-\dy) {receiver actor};
  \draw[pd focus arrow] (\xzero,0)--(\xe,0);
  \node[pd focus datum] at (\sendx,0) {};
  \node[pd label,anchor=east] at ({\sendx-.15},.48) {send};
  \node[pd badge] at ({\sendx+.22},.48) {1};
  \node[pd artifact,minimum width=2.6cm] (MB) at (5.85,-.6) {mailbox};
  \draw[pd thin arrow,rounded corners=2pt] (\sendx,-.08)--(\sendx,-.6)--(MB.west);
  \draw[pd thin arrow,rounded corners=2pt] (MB.east)--(\recvx,-.6)--(\recvx,{-\dy+.08});
  \draw[pd hairline] (\xzero,-\dy)--(\recvx,-\dy);
  \draw[pd focus arrow] (\recvx,-\dy)--(\xe,-\dy);
  \node[pd focus datum] at (\recvx,-\dy) {};
  \node[pd label,anchor=east] at ({\recvx-.15},{-\dy-.48}) {receive};
  \node[pd badge] at ({\recvx+.22},{-\dy-.48}) {2};
  \draw[pd thin arrow] (\xzero,{-\dy-.95})--(\xe,{-\dy-.95});
  \node[pd label,anchor=north east] at (\xe,{-\dy-1.05}) {time};
  \node[pd title,anchor=west] at (0,{\py+.7}) {Synchronous CSP rendezvous};
  \node[pd label,anchor=east] at (1.75,\py) {process $P_1$};
  \node[pd label,anchor=east] at (1.75,{\py-\dy}) {process $P_2$};
  \def\arrSx{4.0}\def\arrRx{5.35}\def\rvx{6.65}
  \foreach \y/\arr in {\py/\arrSx,{\py-\dy}/\arrRx}{
    \draw[pd protocol rule] (\xzero,\y)--(\arr,\y);
    \draw[pd hairline] (\arr,\y)--(\rvx,\y);
    \draw[pd protocol rule] (\rvx,\y)--(\xe,\y);}
  \draw[pd protocol rule] (\rvx,\py)--(\rvx,{\py-\dy});
  \node[pd protocol datum] at (\rvx,{\py-.5*\dy}) {};
  \node[pd badge] at ({\rvx-.55},{\py-.5*\dy}) {3};
  \node[pd protocol label,anchor=west,text width=2.85cm,align=left]
    at ({\rvx+.2},{\py-.5*\dy}) {joint communication};
  \node[pd breach state,minimum width=2.65cm] (TO) at (8.1,{\py-\dy-1.1}) {lease timeout};
  \draw[pd breach arrow,rounded corners=2pt]
    (\arrSx,\py)--(\arrSx,{\py-\dy-1.1})--(TO.west);
  \node[pd badge] at (4.55,{\py-\dy-1.45}) {4};
  \draw[pd thin arrow] (\xzero,{\py-\dy-1.95})--(\xe,{\py-\dy-1.95});
  \node[pd label,anchor=north east] at (\xe,{\py-\dy-2.05}) {time};
\end{tikzpicture}
\par\smallskip
{\pdfiglabel\setlength{\tabcolsep}{3pt}%
\begin{tabularx}{\linewidth}{@{}lX@{}}
\toprule\textbf{Event} & \textbf{Semantics} \\\midrule
1 Actor send & Sender proceeds without receiver readiness. \\
2 Actor receive & Receiver drains the mailbox and may adopt a new behaviour. \\
3 CSP rendezvous & One abstract joint event; the earlier arrival waits. \\
4 Timeout choice & An optional modeled guard offers another transition. \\\bottomrule
\end{tabularx}}
\par\smallskip{\pdfiglabelitalic Delays are illustrative. Ordering depends on the actor
implementation; rendezvous does not imply zero physical communication latency.\par}
\end{minipage}
